\subsection{What Is a Language Model?}
\label{sec:what-is-a-language-model}
\slidesrc{Language models \& RNNs: 2--11}

\subsubsection{Predicting and evaluating language}
\slidesrc{Language models \& RNNs: 3--6}

A \textbf{language model} is a mathematical model that describes human language. It
should let us do two things:
\begin{itemize}[nosep]
  \item \textbf{predict} language: continue a text that has been started;
  \item \textbf{evaluate} language: say how ``natural'' or likely a piece of text is.
\end{itemize}

How can we build such a model? Two facts make it hard to write it down as a fixed set of
rules:
\begin{itemize}[nosep]
  \item \textbf{language is ambiguous}: the same beginning can be continued in many
        reasonable ways (``she went to \ldots'' the shop? work? bed?);
  \item \textbf{language is evolving}: new words and new ways of saying things appear
        all the time, so any fixed list of rules goes out of date.
\end{itemize}
Both facts push in the same direction: instead of saying \emph{which} continuation is
correct, we say \emph{how likely} each continuation is. We therefore use
\textbf{probabilistic models of language}.

\subsubsection{Next-word prediction}
\slidesrc{Language models \& RNNs: 7--9}

The fundamental goal of a language model is \textbf{prediction}, and the simplest form
of prediction is guessing the \textbf{next word}. We number the words of a sentence by
their position: $\vect x^{(1)}$ is the first word, $\vect x^{(2)}$ the second, and so on
(the superscript $(t)$ is a position, ``time step $t$'', not a power).

\begin{center}
\begin{tikzpicture}[word/.style={font=\small, anchor=base, inner sep=2pt},
                    lab/.style={font=\scriptsize, nbgray}]
  \foreach \w/\x [count=\i] in {Anna/0, was/1.25, tired/2.45, so/3.55, she/4.5,
                                went/5.5, to/6.55} {
    \node[word] (w\i) at (\x,0) {\w};
    \node[lab] at (\x,-0.5) {$\vect x^{(\i)}$};
  }
  \node[word, draw=nbred, dashed, rounded corners, minimum width=1.1cm] (w8) at (7.85,0) {\strut ?};
  \node[lab, nbred] at (7.85,-0.55) {$\vect x^{(8)}$};
\end{tikzpicture}
\captionof{figure}{Next-word prediction: having read the first seven words, the model
must say which word is likely to come in position 8.}
\end{center}

Given the words $\vect x^{(1)}, \vect x^{(2)}, \ldots, \vect x^{(t)}$ read so far, the
model computes the \textbf{probability distribution of the next word}
$\vect x^{(t+1)}$:
\[
  p\bigl(\vect x^{(t+1)} \mid \vect x^{(t)}, \ldots, \vect x^{(1)}\bigr).
\]

\notinslides
A few words on this notation, since it is used everywhere in this section.
\begin{itemize}[nosep]
  \item The \textbf{vocabulary} $V$ is the set of all words the model knows, and $|V|$
        is its size (tens of thousands of words for a real model).
  \item A \textbf{probability distribution over the next word} is a list of $|V|$
        numbers, one per word of the vocabulary, each between 0 and 1, which add up to
        1. For example, after ``she went to'' a model could give \emph{the}: 0.5,
        \emph{work}: 0.25, \emph{school}: 0.2, \emph{bed}: 0.05, and 0 to all other
        words.
  \item The vertical bar ``$\mid$'' reads ``given'':
        $p(A \mid B)$ is the probability of $A$ \emph{knowing that} $B$ happened.
        What is written after the bar is called the \textbf{context} or the
        \textbf{conditioning}.
\end{itemize}

Language models are part of daily life: the slide shows a search engine completing a
query while it is being typed (each suggestion is a likely continuation of the
characters typed so far) and a chat assistant writing an answer word by word.

\subsubsection{Evaluating a text: the chain rule}
\slidesrc{Language models \& RNNs: 10}

An equivalent goal is to \textbf{evaluate} a text, i.e.\ assign it a probability. Given
the words $\vect x^{(1)}, \ldots, \vect x^{(T)}$ ($T$ is the length of the text), we want
$p(\vect x^{(1)}, \ldots, \vect x^{(T)})$: the probability that a text is exactly this
sequence of words. Why is this ``equivalent'' to next-word prediction? Because of the
\textbf{chain rule of probability}.

The starting point is the definition of conditional probability:
\[
  p(A, B) = p(A \mid B)\, p(B),
\]
``the probability that both $A$ and $B$ happen is the probability of $B$, times the
probability of $A$ once we know $B$ happened''. Applying it again and again we can peel
a sentence one word at a time, from the end. For three words:
\begin{align*}
  p(\vect x^{(1)}, \vect x^{(2)}, \vect x^{(3)})
  &= p(\vect x^{(3)} \mid \vect x^{(2)}, \vect x^{(1)})\; p(\vect x^{(1)}, \vect x^{(2)})
     \reason{$A = \vect x^{(3)}$, $B = (\vect x^{(1)}, \vect x^{(2)})$} \\
  &= p(\vect x^{(3)} \mid \vect x^{(2)}, \vect x^{(1)})\; p(\vect x^{(2)} \mid \vect x^{(1)})\;
     p(\vect x^{(1)}) \reason{again, on $p(\vect x^{(1)}, \vect x^{(2)})$}
\end{align*}
Reading the last line from right to left: probability of the first word, times the
probability of the second word given the first, times the probability of the third
given the first two. In general:

\begin{nbfact}{Chain rule for a text}
\begin{align*}
  p(\vect x^{(1)}, \ldots, \vect x^{(T)})
  &= p(\vect x^{(1)}) \times p(\vect x^{(2)} \mid \vect x^{(1)}) \times \cdots \times
     p(\vect x^{(T)} \mid \vect x^{(T-1)}, \ldots, \vect x^{(1)}) \\
  &= \prod_{t=1}^{T} \underbrace{p(\vect x^{(t)} \mid \vect x^{(t-1)}, \ldots,
     \vect x^{(1)})}_{\text{next-word prediction!}}
\end{align*}
($\prod$ is the product symbol, the multiplicative analogue of $\sum$; for $t = 1$ the
context is empty and the factor is just $p(\vect x^{(1)})$.)
\end{nbfact}

So the two goals are really one: a model that can predict the next word can evaluate
any text (multiply the next-word probabilities of all its words), and a model that can
evaluate texts can predict the next word (compare the probabilities of the possible
continuations). This identity is exact: no approximation has been made yet.

\begin{nbexample}{probability of a short sentence}
\notinslides
Suppose a model says $p(\textit{the}) = 0.06$ for the first word,
$p(\textit{cat} \mid \textit{the}) = 0.01$ and
$p(\textit{sleeps} \mid \textit{the cat}) = 0.05$. Then
\[
  p(\textit{the cat sleeps}) = 0.06 \times 0.01 \times 0.05 = 0.00003 .
\]
The number is tiny, as is the probability of any particular sentence (there are
astronomically many possible sentences), but it can be \emph{compared}: a good model
gives \emph{the cat sleeps} a much higher probability than \emph{sleeps cat the}.
\end{nbexample}

\subsubsection{Why care about language models?}
\slidesrc{Language models \& RNNs: 11}

\begin{nbdef}{Language model}
A \textbf{language model} is a probabilistic model of natural language: it assigns a
probability to every sequence of words, or equivalently gives, for every context, the
probability distribution of the next word.
\end{nbdef}

\begin{itemize}[nosep]
  \item Language modelling is a good \textbf{benchmark} for general NLP models: to
        predict the next word well a model must know grammar, meaning and facts about
        the world.
  \item Every task that needs to \textbf{estimate the probability of a text} can be
        built on top of an LM: speech recognition, handwriting recognition, spell
        checking, machine translation, summarisation, dialogue, \ldots
  \item Today \textbf{everything else in NLP is built on top of language modelling}:
        GPT is a language model.
\end{itemize}

\begin{nbintuition}[how a probability of text helps other tasks]
\notinslides
A speech recogniser hears a sound that could be ``recognise speech'' or ``wreck a nice
beach''. Both match the audio; the language model says the first is a far more likely
English sentence, so it wins. A spell checker that sees ``I went to the sea \emph{shor}''
proposes \emph{shore} rather than \emph{short} because $p(\textit{shore} \mid
\textit{the sea})$ is much higher. In both cases the LM provides the knowledge of
``what people usually say''.
\end{nbintuition}
